\subsection{伽罗瓦覆叠}

定义 3. —— 设 B 为非空拓扑空间。称 B 的覆叠 E 为伽罗瓦覆叠，如果它是连通的，且对 B 的每个点 b，由 E 的 B-自同构组成的群 \(\mathrm{Aut}_{B}(E)\) 在纤维 \(E_{b}\) 上的作用是传递的。

设 E 为拓扑空间 B 的伽罗瓦覆叠，p 为其投影。映射 p 是满射（A, I, 第 56 页, 定义 6），故空间 B 是连通的。因而覆叠 (E, p) 具有次数，且该次数非零。

命题 5. —— 设 B 为非空拓扑空间，E 为 B 的伽罗瓦覆叠。映射 \((h,x)\mapsto h(x)\) （从 \(\mathrm{Aut}_{\mathrm{B}}(\mathrm{E})\times\mathrm{E}\) 到 E 中）是群 \(\mathrm{Aut}_{\mathrm{B}}(\mathrm{E})^{\circ}\) 在 E 上的一个右作用，它使 E 成为以 \(\mathrm{Aut}_{\mathrm{B}}(\mathrm{E})^{\circ}\) 为群的主覆叠。

给群 \(\mathrm{Aut}_{\mathrm{B}}(\mathrm{E})^{\circ}\) 赋予离散拓扑；此时运算法则 \((h,x)\mapsto h(x)\) 是连续的。此作用是自由的（I, 第 34 页, 命题 11 的推论 2）。由于 E 是 B 的伽罗瓦覆叠，其纤维都是该作用的轨道。由 I, 第 97 页命题 4，E 是以 \(\mathrm{Aut}_{\mathrm{B}}(\mathrm{E})^{\circ}\) 为群的主覆叠。

定理 2. —— 设 B 为连通拓扑空间，E 为 B 的连通非空覆叠。下列性质等价：

\begin{enumerate}
\def\labelenumi{(\roman{enumi})}
\item
  覆叠 E 是伽罗瓦的。
\item
  存在离散拓扑群 G 以及 G 在 E 上的连续右作用，使 E 成为以 G 为群的主覆叠。
\item
  E 的覆叠 \(E \times_{B} E\) （由投影 \(\mathrm{pr}_{1}\) 所定义）是可平凡化的。
\end{enumerate}

当这些条件成立时，由 G 的作用所定义的典范映射 \(G \to \mathrm{Aut}_{B}(E)^{\circ}\) 是群同构。

此外，若空间 B 是局部连通的，则上述性质等价于下述性质：

(i') 存在 B 的点 b，使得群 \(\mathrm{Aut}_{\mathrm{B}}(\mathrm{E})\) 在纤维 \(E_{b}\) 上的作用是传递的。

蕴涵 (i)⇒(ii) 由命题 5 得出，蕴涵 (ii)⇒(iii) 由 I, 第 95 页注 1 得出。我们来证明 (iii)⇒(i)。以 \(p\colon E \to B\) 表示覆叠 E 的投影。由于 B 连通，该覆叠具有次数，又因 E 非空，此次数非零。设 \(b\) 为 B 的一点。我们来证明 \(\mathrm{Aut}_{B}(E)\) 在纤维 \(E_{b}\) 上的作用是传递的。由前所述，该纤维非空。设 \(x\) 与 \(x'\) 为 \(E_{b}\) 中的点。B-态射 \(h\colon E \to E\) （满足 \(h(x) = x'\) ）与连续截面 \(s\) 成一一对应，后者是映射 \(\mathrm{pr}_{1}\colon E \times_{B} E \to E\) 的截面，且满足 \(s(x) = (x, x')\) （I, 第 9 页, 命题 3）。在假设 (iii) 之下，这样的截面存在（I, 第 70 页, 推论 2），且因空间 E 连通而是唯一的（I, 第 34 页, 命题 11 的推论 1）。于是，对每个偶 \((x, x') \in E \times_{B} E\)，存在唯一的 B-态射 \(h\colon E \to E\) 使得 \(h(x) = x'\) 。若 \(h'\colon E \to E\) 是满足 \(h'(x') = x\) 的唯一 B-态射，则有 \(h'(h(x)) = x\) 与 \(h(h'(x')) = x'\) ，从而 \(h' \circ h = \text{Id}_E\) 且 \(h \circ h' = \text{Id}_E\) 。这就证明 \(h\) 是 E 的 B-自同构。因而覆叠 E 是伽罗瓦的。

我们已证明条件 (i)、(ii)、(iii) 等价。设它们成立。以 \(\delta\colon\mathbf{G}\to\mathrm{Aut}_{\mathrm{B}}(\mathrm{E})\) 表示把 \(g\in\mathrm{G}\) 对应到 E 的 B-自同构 \(x\mapsto x\cdot g\) 的映射。映射 \(\delta\) 是从 G 到 \(\mathrm{Aut}_{\mathrm{B}}(\mathrm{E})^{\circ}\) 中的群同态。由于 G 自由地作用于 E，映射 \(\delta\) 是单射。设 \(h\in\mathrm{Aut}_{\mathrm{B}}(\mathrm{E})\)，\(x\in\mathrm{E}\) ，并设 \(g\) 为 G 中满足 \(x\cdot g=h(x)\) 的唯一元素。

B-态射 \(h\) 与 \(\delta(g)\) 在点 \(x\) 处取值相同，从而由于空间 E 连通而处处相同（I, 第 34 页, 命题 11 的推论 1），这就证明了 \(\delta\) 是同构。

现在设空间 B 局部连通，我们证明在假设 (i') 之下覆叠 E 是伽罗瓦的。设 \(E'\) 为 E 按群 \(\mathrm{Aut}_{B}(E)^{\circ}\) 的右作用所得的商空间，并以 \(q\colon E \to E'\) 表示典范映射。它使 E 成为 E' 的满覆叠（I, 第 98 页, 例 2）；映射 \(p\colon E \to B\) 过渡到商，定义出连续映射 \(p'\colon E' \to B\) ，满足 \(p'\circ q = p\) 。由于空间 B 局部连通，B-空间 \((E', p')\) 是一个覆叠（I, 第 81 页, 命题 7）。在假设 (i') 之下，存在 B 的点 \(b\) 使得纤维 \(E'_{b}\) 恰有一个元素；由于空间 B 连通，对 B 的每个点 \(b\) 同样成立（I, 第 74 页, 命题 4），这就证明了 E 是 B 的伽罗瓦覆叠。

命题 6. —— 设 B 为拓扑空间，G 为群。设 (E, p) 为 B 的以 G 为群的主覆叠。假设空间 B 连通且局部连通。设 E₀ 为 E 的一个连通分支，G₀ 为 G 中稳定 E₀ 的子群（A, I, 第 51 页）。B-空间 (E₀, p \textbar{} E₀) 对于 G₀ 在 E₀ 上的右动作是一个主覆叠。此覆叠是伽罗瓦的。

由于空间 B 局部连通，E 亦然，故 \(E_{0}\) 是 E 的开子空间（TG, I, 第 85 页, 命题 11）。又因它还是闭的（TG, I, 第 83 页），空间 \(E_{0}\) 因而是 B 的覆叠（I, 第 80 页, 推论 1）。由于 B 连通，该覆叠具有次数；由于 \(E_{0}\) 非空，\(p|E_{0}\) 的一切纤维都非空。设 x 为 \(E_{0}\) 的一点，并设 \(x' \in E_{0}\) 满足 \(p(x') = p(x)\)；存在元素 \(g \in G\) 使得 \(x \cdot g = x'\)。由于连通分支 \(E_{0}\) 与 \(E_{0} \cdot g\) 有公共点，二者相等，由此得 \(g \in G_{0}\)。于是，群 \(G_{0}\) 传递地作用于 B-空间 \(E_{0}\) 在 \(p(x)\) 处的纤维上。命题 6 得证。

注记. —— 若在命题 6 中不假设空间 B 局部连通，则空间 \(E_{0}\) 可能不是 B 的覆叠（I, 第 147 页, 习题 1）。

